% Referanseliste: alle kilder sitert i decken.
% Verifisert mot paper/references.bib, nedlastede rapporter i
% literature/llm-selskapsrapporter/ og nettsøk (Korinek, Narayanan & Kapoor).

\begin{frame}{Referanser (1/2): akademiske arbeider}
\scriptsize
\begin{columns}[T,onlytextwidth]
\begin{column}{0.48\textwidth}
\begin{itemize}\setlength{\itemsep}{2pt}
  \item Acemoglu, D.\ og D.\ Autor (2011): Skills, Tasks and Technologies. \emph{Handbook of Labor Economics}.
  \item Acemoglu, D.\ og P.\ Restrepo (2019): Automation and New Tasks: How Technology Displaces and Reinstates Labor. \emph{Journal of Economic Perspectives}.
  \item Autor, D., F.\ Levy og R.\ Murnane (2003): The Skill Content of Recent Technological Change. \emph{Quarterly Journal of Economics}.
  \item Bhuller, M., K.\,O.\ Moene, M.\ Mogstad og O.\ Vestad (2022): Facts and Fantasies about Wage Setting and Collective Bargaining. \emph{Journal of Economic Perspectives}.
  \item Brynjolfsson, E., B.\ Chandar og R.\ Chen (2025): Canaries in the Coal Mine? Six Facts about the Recent Employment Effects of AI. Arbeidsnotat, Stanford Digital Economy Lab.
  \item Demirer, M., J.\ Horton, N.\ Immorlica og B.\ Lucier (2026): Chaining Tasks, Redefining Work: A Theory of AI Automation. Arbeidsnotat.
  \item Eloundou, T., S.\ Manning, P.\ Mishkin og D.\ Rock (2024): GPTs are GPTs: Labor Market Impact Potential of LLMs. \emph{Science}.
  \item Felten, E., M.\ Raj og R.\ Seamans (2021): Occupational, Industry, and Geographic Exposure to AI. \emph{Strategic Management Journal}.
  \item Handa, K.\ mfl.\ (2025): Which Economic Tasks are Performed with AI? Evidence from Millions of Claude Conversations. Anthropic/arXiv.
  \item Hernæs, Ø.\ og A.\,R.\ Kostøl (2026): KI-indeksen: Kunstig intelligens og det norske arbeidsmarkedet. kiindeksen.no.
\end{itemize}
\end{column}
\begin{column}{0.48\textwidth}
\begin{itemize}\setlength{\itemsep}{2pt}
  \item Humlum, A.\ og E.\ Vestergaard (2026): Still Waters, Rapid Currents: Early Labor Market Transformation under Generative AI. NBER.
  \item Kauhanen, A.\ og P.\ Rouvinen (2026): AI Has Not Impacted the Youth Labor Market in Finland. ETLA.
  \item Klein Teeselink, B.\ (2025): Generative AI and Labor Market Outcomes: Evidence from the United Kingdom. Arbeidsnotat.
  \item Klein Teeselink, B.\ og D.\ Carey (2026): AI, Automation, and Expertise. Arbeidsnotat.
  \item Korinek, A.\ (2025): AI Agents for Economic Research. NBER Working Paper 34202.
  \item Lindenlaub, I., R.\ Oh, M.\,A.\ Rodríguez og L.\ Veldkamp (2026): Beyond Exposure: Predicting AI Adoption Based on Comparative Advantage. Arbeidsnotat.
  \item Lodefalk, M., L.\ Löthman, M.\ Koch og E.\ Engberg (2026): Same Storm, Different Boats: Generative AI and the Age Gradient in Hiring. Arbeidsnotat.
  \item Narayanan, A.\ og S.\ Kapoor (2025): AI as Normal Technology. Knight First Amendment Institute.
  \item Rambachan, A.\ og J.\ Roth (2023): A More Credible Approach to Parallel Trends. \emph{Review of Economic Studies}.
  \item Webb, M.\ (2020): The Impact of Artificial Intelligence on the Labor Market. Arbeidsnotat.
  \item Yagan, D.\ (2019): Employment Hysteresis from the Great Recession. \emph{Journal of Political Economy}.
\end{itemize}
\end{column}
\end{columns}
\end{frame}

% ---------------------------------------------------------------
\begin{frame}{Referanser (2/2): rapporter og datakilder}
\scriptsize
\begin{columns}[T,onlytextwidth]
\begin{column}{0.48\textwidth}
\begin{itemize}\setlength{\itemsep}{2pt}
  \item Anthropic (2025a): Anthropic Economic Index: AI's Impact on Software Development.
  \item Anthropic (2025b): Anthropic Economic Index: Uneven Geographic and Enterprise AI Adoption (Appel mfl., september 2025).
  \item Anthropic (2026a): Anthropic Economic Index: Economic Primitives (Appel, Massenkoff og McCrory, januar 2026).
  \item Anthropic (2026b): Labor Market Impacts of AI: A New Measure and Early Evidence (Massenkoff og McCrory, mars 2026).
  \item Google (2026): AI \& Economy ATLAS v1.0: Mapping Gemini Usage in the Economy. Google AI \& Economy Research Program, juli 2026.
  \item METR (2025): Measuring AI Ability to Complete Long Tasks.
  \item Microsoft (2026): Global AI Diffusion: Q1 2026 Trends and Insights. Microsoft AI Economy Institute.
\end{itemize}
\end{column}
\begin{column}{0.48\textwidth}
\begin{itemize}\setlength{\itemsep}{2pt}
  \item OpenAI (2026): The AI Jobs Transition Framework: Mapping AI's Near-Term Impact on Jobs. April 2026.
  \item SSB (2011): Standard for yrkesklassifisering (STYRK-08). Notater 17/2011.
  \item SSB (2025): Bruken av KI har skutt fart det siste året (Walther-Zhang og Rybalka); tabell 13265, Bruk av IKT i næringslivet.
  \item U.S.\ Bureau of Labor Statistics (2012, 2017): ISCO-08--SOC- og SOC 2010--2018-krysskoblinger.
  \item O*NET Resource Center, U.S.\ Department of Labor: O*NET-databasen.
  \item A-ordningen via microdata.no (SSB/Sikt/NSD): registerdata for sysselsetting.
\end{itemize}
\end{column}
\end{columns}
\end{frame}
